\chapter{Diagnostika predstavitev modela}
\label{app:diagnostika}

\TODO{
\begin{itemize}
    \item Dodaj variance po plasteh, kosinusne podobnosti, porazdelitve logitov in dodatne slike kolapsa predstavitev.
    \item Dodaj krivulje učne in validacijske izgube po posameznih testnih množicah, če so v glavnem besedilu prikazane samo agregirane ali reprezentativne krivulje.
    \item Dodaj porazdelitve entropije napovednih verjetnosti in razpona logitov za osnovni \ModelName{GCN} ter predlagano GNN.
    \item Dodaj porazdelitve naučenih uteži povezav po relacijah \texttt{temporal\_forward}, \texttt{temporal\_backward}, \texttt{spatial} in \texttt{fixation}, če so artefakti na voljo.
    \item UMAP naučenih predstavitev grafa dodaj sem, če je slika uporabna kot diagnostika, vendar ni dovolj močna za glavno besedilo. Pri vsakem UMAP prikazu zapiši, ali je barvan po razredu, udeležencu ali posnetku.
    \item Če UMAP kaže predvsem subjektne ali stimulusne gruče, ga interpretiraj kot omejitev oziroma opozorilo, ne kot dokaz razredne ločljivosti.
\end{itemize}
}

\section{Diagnostika prekomernega glajenja}
\label{app:diagnostika-prekomerno-glajenje}

\TODO{V tem razdelku poročaj diagnostiko iz
\texttt{GFM-for-eyetracker/results/quick\_v1\_v2\_comparison/RETAIN\_2026-06-12\_16-29-08}.
Interpretacija naj bo informativna: rezultati ne kažejo jasnega kolapsa
predstavitev oziroma prekomernega glajenja pri glavnih grafovskih modelih.
Posebej omeni, da je signalni nabor \texttt{pupil\_only} šibkejši in ima bolj
neodločne napovedi, vendar tega ne predstavimo kot dokaz sistemskega kolapsa
GNN.}

% Vir podatkov: GFM-for-eyetracker/results/quick_v1_v2_comparison/RETAIN_2026-06-12_16-29-08/tables/training_diagnostics_summary.csv.
\begin{table}[H]
    \centering
    \caption{Diagnostika naučenih predstavitev in izhodov najboljšega grafovskega modela po množicah signalov. Makro F1 in uravnotežena točnost opisujeta uspešnost klasifikacije, razpon logitov in entropija gotovost napovedi, varianca predstavitev razpršenost naučenih vektorjev, kos. pod. pa povprečno kosinusno podobnost med predstavitvami.}
    \label{tab:diagnostika-prekomerno-glajenje}
    \footnotesize
    \setlength{\tabcolsep}{2pt}
    \begin{tabularx}{\textwidth}{@{}>{\raggedright\arraybackslash}p{0.13\textwidth} >{\raggedright\arraybackslash}p{0.16\textwidth} *{6}{>{\centering\arraybackslash}X}@{}}
        \toprule
        \textbf{Množica signalov} & \textbf{Najboljši GNN} & \makecell{\textbf{Makro}\\\textbf{F1}} & \makecell{\textbf{Urav.}\\\textbf{točnost}} & \makecell{\textbf{Razpon}\\\textbf{logitov}} & \makecell{\textbf{Entro-}\\\textbf{pija}} & \makecell{\textbf{Varianca}\\\textbf{predst.}} & \makecell{\textbf{Kos.}\\\textbf{pod.}} \\
        \midrule
        vsi signali & \makecell[l]{\ModelName{HeteroGCN-}\\\ModelName{MLP-w}} & $0{,}608$ & $0{,}623$ & $3{,}31$ & $0{,}545$ & $0{,}352$ & $0{,}409$ \\
        samo pogled & \makecell[l]{\ModelName{HeteroGCN-}\\\ModelName{MLP}} & $0{,}658$ & $0{,}662$ & $3{,}93$ & $0{,}583$ & $0{,}291$ & $0{,}307$ \\
        pogled in zenici & \makecell[l]{\ModelName{HeteroGCN-}\\\ModelName{MLP-w}} & $0{,}644$ & $0{,}652$ & $3{,}38$ & $0{,}582$ & $0{,}336$ & $0{,}360$ \\
        samo zenici & \ModelName{GCN} & $0{,}584$ & $0{,}607$ & $1{,}40$ & $0{,}653$ & $0{,}258$ & $0{,}569$ \\
        \bottomrule
    \end{tabularx}
\end{table}

